% -----------------------------------------------------------------------------
% MAIN CONTENT — USER'S NEED PRESENTATION
% -----------------------------------------------------------------------------
\section{Main Content}

\subsection{3. User's Need Presentation}

This section characterizes the operational security environment, end-user civil-protection needs, architectural allocation, and external interfaces that govern the \abbr{EWAS} Space Segment (\texttt{CI-03}, \abbr{WBS} Element \texttt{123456.03}) as the primary sensing and space-relay component of the overall \abbr{EWAS} system (\textbf{AD-01}).

\subsubsection{3.1 Operational Context, Threat Landscape, and Mission Need}

The \textbf{Baltic Sea region} ($53^\circ\text{N}\text{--}66^\circ\text{N}$, $10^\circ\text{E}\text{--}32^\circ\text{E}$) faces an increasingly volatile strategic security environment. A critical vulnerability of the region is its close geographic proximity to potential launch corridors for Short-Range Ballistic Missiles (\abbr{SRBM}, range $< 1{,}000\text{ km}$, flight time 3--7~min) and Medium-Range Ballistic Missiles (\abbr{MRBM}, range 1,000--3,000~km, flight time 7--15~min). The region is also exposed to Intermediate-Range and Intercontinental Ballistic Missiles (\abbr{IRBM}/\abbr{ICBM}) launched from outside the immediate Baltic theatre whose trajectories threaten Baltic states (\texttt{R-049}).

Conventional early warning postures relying exclusively on ground-based radar installations suffer from fundamental physical and geometric limitations in this theatre:
\begin{itemize}
    \item \textbf{Earth-Curvature Horizon Masking:} Due to Earth's curvature, a terrestrial radar with an antenna height of $h_r \approx 30\text{ m}$ cannot achieve line-of-sight visibility to an ascending missile until the vehicle clears the local radar horizon $d_h \approx 4.12\bigl(\sqrt{h_r} + \sqrt{h_t}\bigr)$ (in km, for altitudes in meters). At typical regional stand-off distances of 400--800~km, an ascending missile remains masked below the radar horizon throughout its initial boost phase (altitudes below 15--45~km). This delays initial detection by 60--150~seconds after rocket motor ignition.
    \item \textbf{Terrain Obstruction, Atmospheric Clutter, and Depressed Trajectories:} Forested topography, maritime ducting, and quasi-ballistic depressed trajectories further degrade low-elevation ground radar acquisition. Consequently, national civil-protection authorities are left with insufficient lead time to execute public sheltering and infrastructure protection protocols.
\end{itemize}

To eliminate this horizon-masking blind spot, end users---specifically \textbf{Authorized National Crisis-Management Authorities, Ministries of the Interior, and Emergency Warning Centers} across the Baltic partner states (Chapter 8 of \textbf{AD-01})---require a persistent, space-based early warning capability. This system must detect ballistic missile launches at the earliest physically observable stage: \textbf{rocket motor ignition and boost-phase ascent} (\texttt{R-007}, \texttt{R-009}). During the boost phase (typically lasting 45--180~seconds), propellant combustion releases megawatts of radiant intensity in the Short-Wave Infrared (\abbr{SWIR}) and Medium-Wave Infrared (\abbr{MWIR}) spectral bands. This provides a high-contrast radiometric signature observable from orbit above cloud tops.

To translate space-based boost-phase detection into actionable civil-protection lead time, the user need imposes three overarching system drivers on \abbr{EWAS}:
\begin{enumerate}
    \item \textbf{Automated Preliminary Warning within $\le 60\text{ Seconds}$ (\texttt{R-011}, \texttt{R-045}):} The elapsed time from the first Geostationary Earth Orbit (\abbr{GEO}) detection of a candidate launch event to the automated delivery of a preliminary warning product at the authorized national authority interface must not exceed \textbf{60~seconds}. Within this 60-second window, the system cannot rely solely on coarse single-tier \abbr{GEO} coordinates. The warning may be issued \textbf{only after} two selected Low Earth Orbit (\abbr{LEO}) satellites have acquired follow-on stereoscopic observations and the Ground Segment has fused those observations into a refined 3D trajectory and probabilistic impact-area (\textbf{Instantaneous Impact Point --- \abbr{IIP}}) estimate with explicit confidence ellipses (\texttt{R-012}, \texttt{R-013}, \texttt{R-045}).
    \item \textbf{Continuous 24/7 Surveillance and $> 99.9\%$ Availability (\texttt{R-006}, \texttt{R-008}, \texttt{R-010}, \texttt{R-015}):} The Space Segment must maintain uninterrupted temporal and spatial surveillance without coverage gaps over the protected region and surrounding launch corridors. It must achieve a launch detection probability $> 95\%$ (\texttt{R-010}) and support an overall system operational availability $> 99.9\%$ (\texttt{R-015}) across 7-year (\abbr{LEO}) and 15-year (\abbr{GEO}) operational design lifetimes (\texttt{R-016}).
    \item \textbf{Strict Compliance with Non-Offensive Civil-Protection and Data Sovereignty Boundaries (\texttt{R-001}--\texttt{R-004}, \texttt{R-048}):} \abbr{EWAS} is conceived exclusively as a defensive civil-protection decision-support asset:
    \begin{itemize}
        \item \emph{Non-Offensive Mandate (\texttt{R-001}):} The system is strictly prohibited from providing weapon targeting coordinates, fire-control solutions, interceptor mid-course guidance, or direct command of military kinetic response assets.
        \item \emph{Scope Limitation (\texttt{R-002}):} The system is not required to perform definitive missile-variant identification, warhead nuclear/conventional yield estimation, or post-impact battle damage assessment.
        \item \emph{Authorized Institutional Access \& Sovereign Archival (\texttt{R-003}, \texttt{R-004}, \texttt{R-034}--\texttt{R-037}, \texttt{R-047}, \texttt{R-048}):} Warning outputs are delivered exclusively to authorized governmental crisis-management interfaces (\texttt{R-003}), which form the external boundary of \abbr{EWAS}. Downstream public siren and cell-broadcast activation remains the responsibility of national authorities (\texttt{R-048}). All archived telemetry and observation records are retained for at least 5~years (\texttt{R-036}) and restricted strictly to algorithm refinement, radiometric/geometric calibration, performance verification, operator training, and post-event technical analysis (\texttt{R-004}).
    \end{itemize}
\end{enumerate}

\subsubsection{3.2 Role of the Space Segment in the Two-Segment EWAS Warning Chain}

The Space Segment (\abbr{SSEG}, \abbr{WBS} \texttt{123456.03}) operates as one of the \textbf{two integrated system segments} that constitute the \abbr{EWAS} baseline (\textbf{AD-01}, Chapters 3 and 4):
\begin{enumerate}
    \item \textbf{Space Segment (\abbr{SSEG}, \abbr{WBS} \texttt{123456.03}, \texttt{CI-03}):} The space-borne sensing, onboard processing, and inter-satellite relay infrastructure. It comprises the 6 \abbr{GEO} detection satellites (\texttt{SATG}, \abbr{WBS} \texttt{123456.03.02}) at 35,786~km and the 72 \abbr{LEO} tracking satellites (\texttt{SATL}, \abbr{WBS} \texttt{123456.03.01}) at 1,000~km, interconnected by a \textbf{mandatory Space Segment Optical Inter-Satellite Link (\abbr{OISL}) network} (1,550~nm laser \abbr{OISL} mesh across \abbr{LEO} and cross-layer \abbr{GEO}--\abbr{LEO} optical relay trunks).
    \item \textbf{Ground Segment (\abbr{WBS} \texttt{123456.02}):} The terrestrial infrastructure comprising the \emph{Data Archiving System} (\texttt{123456.02.01}), \emph{Command \& Control System} (\texttt{123456.02.02}), \emph{Mission Control Center} (\texttt{123456.02.03}), \emph{Data Processing System} (\texttt{123456.02.04}), \emph{Ground Station Network} (\texttt{123456.02.05}), and \emph{Early Warning System} (\texttt{123456.02.06}). External national crisis-management authorities are warning recipients connected to the Early Warning System interface, not a third system segment (\texttt{R-048}).
\end{enumerate}

To satisfy the $\le 60\text{-second}$ end-to-end latency requirement (\texttt{R-011}) from initial \abbr{GEO} detection to external national authority warning dispatch, the operational warning chain is partitioned into four deterministic steps across the Space and Ground Segments. Table~\ref{tab:warning_chain_latency} defines the functional sequence, segment responsibilities, exchanged data products (aligned with the Ground Segment Data Dictionary in Chapter 12 of \textbf{AD-01}), and binding latency budget allocations.

\clearpage
\begin{table}[H]
\centering
\caption{EWAS End-to-End Warning Chain and Space Segment Latency Allocation (\texttt{R-011})}
\label{tab:warning_chain_latency}
{\footnotesize\setlength{\tabcolsep}{4pt}\renewcommand{\arraystretch}{1.12}%
\arrayrulecolor{TableBorderBlue}%
\rowcolors{2}{white}{TableRowLightBlue}%
\noindent
\begin{tabularx}{\textwidth}{|>{\raggedright\arraybackslash}p{2.1cm}|>{\raggedright\arraybackslash}p{2.2cm}|>{\raggedright\arraybackslash}X|>{\raggedright\arraybackslash}p{2.3cm}|>{\raggedright\arraybackslash}p{1.9cm}|}
    \hline
    \rowcolor{TableHeaderBlue}
    \thead{Chain Step} & \thead{Active Element} & \thead{Required Operational Functions} & \thead{Output Product \& Req.\ Ref.} & \thead{Latency Allocation} \\ \hline
    \textbf{Step 1:} \newline Boost-Phase Detection \& Reporting & \textbf{SSEG (\texttt{SATG})} \newline 6~\abbr{GEO} nodes at 35,786~km & Persistent \abbr{WFOV} \abbr{SWIR}/\abbr{MWIR} staring surveillance; real-time onboard proton despiking, background clutter suppression, and plume event classification; immediate encrypted downlink to Ground Station Network. & \emph{GEO Event Report} \newline (\texttt{R-009}, \texttt{R-019}, \texttt{R-022}, \texttt{R-042}) & \textbf{$\le 10.0\text{ s}$} \newline (from $T_0$ plume threshold) \\ \hline
    \textbf{Step 2:} \newline LEO Selection \& \abbr{ISL} Cueing & \textbf{Ground C2} + \textbf{SSEG \abbr{ISL}} \newline (\abbr{GEO}/\abbr{LEO} relay) & Ground Segment validates \abbr{GEO} report, evaluates 1,000~km Walker Delta geometry, selects two optimal \abbr{LEO} satellites, and uplinks tracking requests via direct ground stations or \textbf{mandatory \abbr{SSEG} \abbr{ISL} relay paths}. & \emph{LEO Tracking Request} \newline (\texttt{R-030}, \texttt{R-043}, \texttt{R-044}) & \textbf{$\le 5.0\text{ s}$} \newline (C2 selection $\le 2\text{ s}$; \abbr{ISL} routing $\le 3\text{ s}$) \\ \hline
    \textbf{Step 3:} \newline Agile \abbr{LOS} Lock, Stereo Tracking \& Downlink & \textbf{SSEG (\texttt{SATL})} \newline 2 cued \abbr{LEO} nodes at 1,000~km + \textbf{\abbr{ISL} Mesh} & Agile optical \abbr{LOS} slewing via two-axis gimbal + \abbr{FSM} to achieve target lock-on in $< 5.0\text{ s}$ with $< 2\text{ arcsec/s}$ stability; synchronized stereoscopic \abbr{NFOV} tracking; \abbr{ISL} exchange of timestamped \abbr{LOS} vectors to compute onboard preliminary 3D trajectory; encrypted downlink of \abbr{LOS} packages + preliminary trajectory via direct link or \abbr{ISL}. & \emph{LEO Observation Package} \& \emph{Onboard Preliminary Trajectory} \newline (\texttt{R-023}, \texttt{R-024}, \texttt{R-033}) & \textbf{$\le 25.0\text{--}30.0\text{ s}$} \newline ($< 5\text{ s}$ lock-on; $15\text{--}18\text{ s}$ stereo dwell/3D fit; $\le 5\text{--}7\text{ s}$ downlink) \\ \hline
    \textbf{Step 4:} \newline Ground Fusion \& Warning & \textbf{Ground Segment} \newline (\texttt{123456.02.04} \& \texttt{123456.02.06}) & Central Data Processing System performs authoritative multi-sensor fusion of downlinked stereoscopic \abbr{LOS} vectors (with covariance metadata) and onboard preliminary trajectory to compute refined 3D trajectory, probabilistic \abbr{IIP} ellipse, and Baltic threat priority; Early Warning System automatically dispatches preliminary warning to national authorities prior to operator review. & \emph{Ground Refined Trajectory}, \emph{IIP Assessment} \& \emph{Preliminary Warning} \newline (\texttt{R-012}, \texttt{R-026}, \texttt{R-027}, \texttt{R-045}, \texttt{R-049}) & \textbf{$\le 10.0\text{--}15.0\text{ s}$} \newline (Ground fusion + automated dispatch) \\ \hline
    \textbf{Total Chain} & \textbf{EWAS System} & \textbf{End-to-End Elapsed Time from Initial \abbr{GEO} Detection ($T_0$) to Authority Interface Delivery} & \textbf{Compliant with \texttt{R-011} \& \texttt{R-045}} & \textbf{$\le 60.0\text{ s}$ Total} \newline (\textbf{$\le 40\text{--}45\text{ s}$ \abbr{SSEG}}) \\ \hline
\end{tabularx}}
\end{table}

As shown in Table~\ref{tab:warning_chain_latency}, the Space Segment executes \textbf{Steps 1 and 3} directly and provides the mandatory space-relay transport backbone for \textbf{Step 2}. This represents a combined Space Segment latency allocation of $\le 40\text{--}45\text{ seconds}$. Two architectural mechanisms within \abbr{SSEG} are essential to closing this timing budget:
\begin{itemize}
    \item \textbf{Mandatory Space Segment \abbr{OISL} Routing (\texttt{R-023}, \texttt{R-044}):} The two best-positioned \abbr{LEO} satellites over a launch corridor outside the Baltic region may not be in simultaneous line-of-sight of a European ground station at the instant of launch. Relying solely on direct ground-station passes would introduce multi-minute store-and-forward delays and violate \texttt{R-011}. The mandatory \abbr{SSEG} laser Optical Inter-Satellite Link (\abbr{OISL}, 1,550~nm) network---combining intra-/inter-plane \abbr{OISL} mesh links across the 72 \abbr{LEO} satellites with cross-layer \abbr{GEO}--\abbr{LEO} optical relay trunks---guarantees continuous sub-3-second command delivery (\texttt{R-044}), sub-second inter-satellite \abbr{LOS} vector exchange for onboard preliminary 3D trajectory computation (\texttt{R-023}), and immediate telemetry and observation return.
    \item \textbf{Independent Optical Payload \abbr{LOS} Steering Without Affecting Spacecraft Bus Attitude or Other Instrument Pointing Stability (\texttt{R-024}):} Slewing an entire 1,000~km \abbr{LEO} spacecraft bus through $20^\circ\text{--}45^\circ$ using reaction wheels alone would require 15--30~seconds of settling time and would break the pointing lock of the laser \abbr{OISL} terminals and ground antennas. Instead, each \texttt{SATL} node steers its infrared payload Line of Sight (\abbr{LOS}) independently via a two-axis payload gimbal and high-bandwidth Fast Steering Mirror (\abbr{FSM}) while the spacecraft bus remains 3-axis nadir-stabilized. This enables target lock-on within $< 5\text{ seconds}$ and $< 2\text{ arcsec/s}$ \abbr{LOS} stability (\texttt{R-024}) without disturbing the attitude of the spacecraft bus or the pointing stability of the four laser \abbr{OISL} terminals and space-to-ground antennas.
\end{itemize}

\subsubsection{3.3 Space Segment System Boundaries and External Interfaces}

The boundary of Configuration Item \texttt{CI-03} (\abbr{SSEG}) is defined by four external interface categories connecting the flight constellation to the Ground Segment, the Launch Segment, the observable target/background environment, and regulatory/space-traffic entities (under the project interface-management rules and applicable ICDs). Table~\ref{tab:sseg_external_interfaces} summarizes the external interfaces of the Space Segment.

\begin{table}[H]
\centering
\caption{Summary of Space Segment (\texttt{CI-03}) External Interfaces}
\label{tab:sseg_external_interfaces}
{\footnotesize\setlength{\tabcolsep}{4pt}\renewcommand{\arraystretch}{1.12}%
\arrayrulecolor{TableBorderBlue}%
\rowcolors{2}{white}{TableRowLightBlue}%
\noindent
\begin{tabularx}{\textwidth}{|>{\raggedright\arraybackslash}p{1.9cm}|>{\raggedright\arraybackslash}p{3.1cm}|>{\raggedright\arraybackslash}p{2.9cm}|>{\raggedright\arraybackslash}X|}
    \hline
    \rowcolor{TableHeaderBlue}
    \thead{Interface ID} & \thead{Interface Name} & \thead{External Counterpart} & \thead{Principal Functional, Physical \& Protocol Characteristics} \\ \hline
    \textbf{IF-EXT-01} & Space-to-Ground \abbr{RF} \abbr{TT\&C} \& Mission Data Interface & Ground Station Network (\abbr{WBS} \texttt{123456.02.05}) \& Ground C2 (\texttt{123456.02.02}) & All-weather bidirectional S-band/X-band \abbr{TT\&C} and high-rate X-band/Ka-band mission data downlinks; \abbr{CCSDS} packet/frame compliant; \abbr{COMSEC} encryption (\texttt{R-031}) and spread-spectrum anti-jamming protection (\texttt{R-032}, \texttt{R-033}). \\ \hline
    \textbf{IF-EXT-02} & Launch Vehicle \& Dispenser Interface & Launch Service Providers (Heavy-Lift \abbr{GEO} \& Medium/Heavy \abbr{LEO}) & Standard mechanical separation rings for \texttt{SATG}; multi-satellite rideshare dispenser mechanisms deploying 9~\texttt{SATL} spacecraft per plane; umbilical power/data/purge lines; launch fairing acoustic, vibrational, and quasi-static envelopes. \\ \hline
    \textbf{IF-EXT-03} & Target Plume \& Background Optical Interface & Boost/Post-Boost Missile Plumes \& Earth/Solar Background & Radiometric entrance aperture interface in \abbr{SWIR} ($\sim 2.7~\mu\mathrm{m}$ $\mathrm{H_2O}/\mathrm{CO_2}$ band) and \abbr{MWIR} ($\sim 4.2\text{--}4.4~\mu\mathrm{m}$ $\mathrm{CO_2}$ red/blue spike bands); solar glint, Earth albedo, cloud-top thermal clutter, and stray-light rejection (\texttt{R-009}, \texttt{R-020}, \texttt{R-025}). \\ \hline
    \textbf{IF-EXT-04} & Space Traffic Management \& Regulatory Interface & Space Surveillance Networks \& \abbr{ITU} Radiocommunication Sector & Reception of Conjunction Data Messages (CDMs) via Ground C2 for collision avoidance at 1,000~km \abbr{LEO} and 35,786~km \abbr{GEO}; compliance with \abbr{ITU} frequency allocations and the project space-debris disposal requirements (\texttt{R-018}). \\ \hline
\end{tabularx}}
\end{table}

\noindent\emph{Note on External Space-to-Ground RF Links (\texttt{IF-EXT-01}) vs.\ Internal Space-to-Space Optical Links (\abbr{OISL}):} Table~\ref{tab:sseg_external_interfaces} lists only interfaces that cross the external boundary of Configuration Item \texttt{CI-03} (\abbr{SSEG}). External Space-to-Ground links (\texttt{IF-EXT-01}) utilize Radio Frequency (\abbr{RF}: S-band, X-band, and Ka-band) to guarantee all-weather signal penetration through tropospheric cloud cover and precipitation over European ground stations. By contrast, all Space-to-Space Inter-Satellite Links between \texttt{SATL} and \texttt{SATG} operate in vacuum above the atmosphere and utilize 1,550~nm laser Optical Inter-Satellite Links (\abbr{OISL}); because these links are internal to the Space Segment, they are specified in Section~4.2 and Section~7.7.

\begin{enumerate}
    \item \textbf{Space-to-Ground Interface (\texttt{IF-EXT-01}: \abbr{SSEG} $\leftrightarrow$ Ground Segment \texttt{123456.02}):} Governs all all-weather radio-frequency links between the 6 \texttt{SATG} and 72 \texttt{SATL} spacecraft and the Ground Station Network (\texttt{123456.02.05}). It is managed under the \emph{Space Segment Interface Control Document} (\abbr{WBS} \texttt{123456.05.02.04}) and \emph{Data Exchange ICD} (\texttt{123456.05.02.01}). The uplink carries authenticated telecommands, cryptographic key updates, ephemeris uploads, software patches, and \emph{LEO Tracking Requests} (\texttt{R-043}, \texttt{R-044}). The downlink carries spacecraft housekeeping telemetry, \emph{GEO Event Reports} (\texttt{R-042}), \emph{LEO Observation Packages}, and \emph{Onboard Preliminary Trajectory} estimates (\texttt{R-023}). All links enforce end-to-end encryption (\texttt{R-031}), command authentication, and anti-jamming resilience (\texttt{R-032}).
    \item \textbf{Launch Segment Interface (\texttt{IF-EXT-02}: \abbr{SSEG} $\leftrightarrow$ Launch Vehicles):} Defines the mechanical, electrical, thermal, cleanliness, and dynamic load boundaries between the spacecraft Structures and Mechanisms Subsystem (\abbr{STR/MEC}) and commercial launch vehicles. For \texttt{SATG}, the interface supports single or dual-manifest geostationary transfer or direct \abbr{GEO} injection. For \texttt{SATL}, the interface uses modular multi-satellite dispenser rings that batch-inject a complete orbital plane of 9 \abbr{LEO} satellites per launch into the 1,000~km, $78^\circ\text{--}82^\circ$ inclination orbit across 8 dedicated plane launches.
    \item \textbf{Target Plume and Background Optical Interface (\texttt{IF-EXT-03}: \abbr{SSEG} Payload $\leftrightarrow$ Radiometric Scene):} Defines the uncooperative optical boundary between the \abbr{SWIR}/\abbr{MWIR} telescope entrance pupils and the observed scene. The payload discriminates high-temperature booster exhaust plumes ($\mathrm{H_2O}$ and $\mathrm{CO_2}$ emission around $2.7~\mu\mathrm{m}$ in \abbr{SWIR} and $4.2\text{--}4.4~\mu\mathrm{m}$ in \abbr{MWIR}) against terrestrial and space backgrounds. These backgrounds include solar reflections off cirrus clouds, sunlit limb scatter, industrial or wildfire thermal sources, and South Atlantic Anomaly (\abbr{SAA}) or polar proton transients (\texttt{R-019}--\texttt{R-022}, \texttt{R-025}).
    \item \textbf{Space Traffic Management and Regulatory Interface (\texttt{IF-EXT-04}):} Governs operational coordination with international space situational awareness providers via the Ground Segment Mission Control Center. It covers the ingestion of Conjunction Data Messages (CDMs) for collision-avoidance manoeuvres in the 1,000~km \abbr{LEO} and 35,786~km \abbr{GEO} regimes, adherence to International Telecommunication Union (\abbr{ITU}) spectral masks, laser aviation-safety exclusions for optical links, and the project space-debris mitigation policy and \textbf{ESA Zero Debris} end-of-life disposal mandate (\texttt{R-018}).
\end{enumerate}
